\chapter{Introduction}
\label{ch:intro}

\section{Context and Motivation}

The electrocardiogram (ECG) is one of the most common heart tests, but interpretation is difficult: clinicians can disagree about the same trace \cite{dhutia_interrater_2017}, and readings need to be reviewed by experienced readers \cite{schlapfer_computer-interpreted_2017}. Many ECGs are stored as prints or photographs, which automated analysis built for digital signals cannot use \cite{shivashankara_ecg-image-kit_2024}.

\section{Problem Statement}
\label{sec:problem}

Deep learning models classify ECGs well when given raw signal data
\cite{hannun_cardiologist_2019,ribeiro_automatic_2020}. Work on images is newer, and
Chapter~\ref{ch:lit} shows that it is limited in four ways:
\begin{enumerate}\setlength{\itemsep}{0pt}\setlength{\parskip}{0pt}
  \item Image-based classifiers have mostly been tested on synthesised images, or on a limited set of real images.
  \item Image-based studies rarely report whether the model's confidence scores are calibrated.
  \item Explanations of the predictions are rarely evaluated.
  \item There is little evidence on how results should be presented to different kinds of users.
\end{enumerate}
A system that flags possible abnormalities from real ECG images, states its confidence and shows what each prediction was based on would address these gaps. It must also cope with rendering shortcuts, a grid that is the only record of scale, and imbalanced classes.

\section{Aim and Objectives}

\textbf{Aim.} To build and evaluate a convolutional neural network (CNN) application that
flags possible abnormalities from ECG images, reports its calibrated confidence, and explains its predictions to different users. The system is a supportive review and is not a diagnostic tool to replace health professionals' advice.

\begin{enumerate}\setlength{\itemsep}{0pt}\setlength{\parskip}{0pt}
  \item \textbf{O1:} Render a labelled training set from PTB-XL with ECG-Image-Kit, and assemble a real test set from the ECG-Image-Database.
  \item \textbf{O2:} Develop a preprocessing and waveform-detection pipeline that detects the trace and the beats.
  \item \textbf{O3:} Train a multi-label CNN and evaluate it with imbalance-aware metrics on patient-wise splits.
  \item \textbf{O4:} Measure the drop in performance from rendered to real images.
  \item \textbf{O5:} Calibrate confidence with temperature scaling; report calibration error per class.
  \item \textbf{O6:} Explain predictions with Grad-CAM in a GUI for different users, and run a user study.
\end{enumerate}